\section{Énumération microcanonique des états sur un écran fini}

Considérons la réalisation finie
\[
 \mathcal H_{\rm occ}=\bigoplus_{k=0}^{3}\CC^{g_k},
 \qquad
 H_{\rm occ}|_{\CC^{g_k}}=E_k^{\rm occ} I,
\]
avec niveaux et multiplicités
\[
 E^{\rm occ}=(0,3,6,9),\qquad g=(1,6,12,8),
 \qquad(N,E)=(12,66).
\]
La somme \(1+6+12+8=27\) fait de \(\mathcal H_{\rm occ}\) un écran à
vingt-sept modes. Ce spectre est la donnée d’une représentation finie
déclarée ; son cardinal ne suffit pas à l’identifier à toute autre
apparition de \(27\) dans HMT\@. Le morphisme qui relie un écran concret
de bord à \((\mathcal H_{\rm occ},H_{\rm occ})\) doit spécifier aussi bien
les fibres que l’opérateur spectral.
Les fonctions génératrices fermionique et bosonique sont
\[
 F(y,x)=\prod_k(1+yx^{E_k^{\rm occ}})^{g_k},
 \qquad
 B(y,x)=\prod_k(1-yx^{E_k^{\rm occ}})^{-g_k}.
\]

\begin{proposition}[Cardinalités microcanoniques]
\label{prop:cardinalidades-microcanonicas}
\[
 [y^{12}x^{66}]F=2\,104\,494,
 \qquad
 [y^{12}x^{66}]B=259\,436\,430.
\]
\end{proposition}
\begin{proof}
Dans le produit fermionique, choisir \(r\) parmi les \(g_k\) modes contribue pour
\(\binom{g_k}{r}y^rx^{rE_k^{\rm occ}}\). Dans le produit bosonique, répartir \(r\)
occupations indiscernables entre \(g_k\) modes contribue pour
\(\binom{g_k+r-1}{r}y^rx^{rE_k^{\rm occ}}\). La convolution finie sur les
quatre niveaux produit les entiers indiqués.
\end{proof}

\section{Distributions de Bose--Einstein et de Fermi--Dirac}

La formulation d’équilibre part du Hamiltonien et du nombre total
d’occupations,
\begin{equation}
 H=\sum_kE_k^{\rm occ} n_k,
 \qquad
 N=\sum_k n_k,
 \qquad
 K_\mu:=H-\mu N.
 \label{eq:hamiltoniano-gran-canonico-rev3}
\end{equation}
Pour \(\beta>0\), l’état grand canonique est
\begin{equation}
 \rho_{\beta,\mu}
 =\frac{e^{-\beta K_\mu}}{\mathcal Z_{\beta,\mu}},
 \qquad
 \mathcal Z_{\beta,\mu}
 =\operatorname{Tr}(e^{-\beta K_\mu}).
 \label{eq:estado-gran-canonico-rev3}
\end{equation}

\begin{theorem}[Factorisation grand canonique et lois d’occupation]
\label{thm:factorizacion-gran-canonica-rev3}
Pour des niveaux d’énergie \(E_k^{\rm occ}\) de dégénérescences \(g_k\), la
trace fermionique se factorise comme
\begin{equation}
 \mathcal Z_{\rm FD}
 =\prod_k\bigl(1+e^{-\beta(E_k^{\rm occ}-\mu)}\bigr)^{g_k},
 \qquad
 \langle n_k\rangle_{\rm FD}
 =\frac{g_k}{e^{\beta(E_k^{\rm occ}-\mu)}+1}.
 \label{eq:factorizacion-fd-rev3}
\end{equation}
Dans le secteur bosonique, sous la condition
\(\mu<\inf_kE_k^{\rm occ}\),
\begin{equation}
 \mathcal Z_{\rm BE}
 =\prod_k\bigl(1-e^{-\beta(E_k^{\rm occ}-\mu)}\bigr)^{-g_k},
 \qquad
 \langle n_k\rangle_{\rm BE}
 =\frac{g_k}{e^{\beta(E_k^{\rm occ}-\mu)}-1}.
 \label{eq:factorizacion-be-rev3}
\end{equation}
\end{theorem}

\begin{proof}
L’opérateur \(K_\mu\) est une somme d’opérateurs d’occupation commutants, de
sorte que sa trace se factorise par modes. Dans chaque mode fermionique, on somme les
occupations \(0\) et \(1\), et le facteur correspondant est
\(1+u_k\), où
\(u_k=e^{-\beta(E_k^{\rm occ}-\mu)}\). Dans chaque mode bosonique, on somme toutes
les occupations non négatives et on obtient la série géométrique
\((1-u_k)^{-1}\) ; la condition sur \(\mu\) garantit \(u_k<1\). Les
puissances \(g_k\) incorporent la dégénérescence. Enfin,
\[
 \langle n_k\rangle
 =-\frac1\beta\frac{\partial}{\partial E_k^{\rm occ}}
   \log\mathcal Z
\]
produit les deux lois d’occupation
\cite{Bose1924,Einstein1924,Fermi1926,Dirac1926,Pathria2011}.
\end{proof}

La dérivation microcanonique retrouve le même résultat dans le régime
thermodynamique. Pour des occupations fermioniques moyennes
\(p_k=m_k/g_k\), l’entropie combinatoire par niveau est
\begin{equation}
 S_k\simeq
 g_k\bigl[-p_k\log p_k-(1-p_k)\log(1-p_k)\bigr].
 \label{eq:entropia-fd-nivel-rev3}
\end{equation}
La variation de \(\sum_kS_k\), soumise à
\(\sum_kg_kp_k=N\) et
\(\sum_kg_kE_k^{\rm occ}p_k=E\), donne
\begin{equation}
 \log\frac{1-p_k}{p_k}=\alpha+\beta E_k^{\rm occ},
 \qquad
 \mu=-\frac{\alpha}{\beta},
 \label{eq:variacion-fd-rev3}
\end{equation}
et, par conséquent, l’occupation de Fermi--Dirac de \eqref{eq:factorizacion-fd-rev3}. Les multiplicateurs \(\beta\) et \(\mu\) sont les coordonnées duales des contraintes d’énergie et de population.

Pour des occupations bosoniques moyennes
\(\overline n_k=m_k/g_k\), la log-multiplicité du niveau est
\begin{equation}
 S_k\simeq
 g_k\bigl[(1+\overline n_k)\log(1+\overline n_k)
          -\overline n_k\log\overline n_k\bigr].
 \label{eq:entropia-be-nivel-rev3}
\end{equation}
La variation de \(\sum_kS_k\), soumise à
\(\sum_kg_k\overline n_k=N\) et
\(\sum_kg_kE_k^{\rm occ}\overline n_k=E\), produit
\begin{equation}
 \log\frac{1+\overline n_k}{\overline n_k}
 =\alpha+\beta E_k^{\rm occ},
 \qquad
 \overline n_k
 =\frac{1}{e^{\beta(E_k^{\rm occ}-\mu)}-1},
 \qquad
 \mu=-\frac{\alpha}{\beta}.
 \label{eq:variacion-be-rev3}
\end{equation}
C’est l’occupation de Bose--Einstein de
\eqref{eq:factorizacion-be-rev3}. La condition
\(\mu<\inf_kE_k^{\rm occ}\) assure la positivité du dénominateur et la
convergence de chaque facteur modal.

Dans l’instance finie construite, les deux contraintes déterminent
\[
 (\beta_{\rm FD},\mu_{\rm FD})
 =(0.15307011353415928,\,4.502865197597043),
\]
\[
 (\beta_{\rm BE},\mu_{\rm BE})
 =(0.05456727960068316,\,-15.92569781906737).
\]
Si \(m_k^{\rm mic}\) désigne l’occupation moyenne obtenue par énumération
microcanonique exacte et \(m_k^{\rm eq}\) l’évaluation de la distribution
correspondante, les erreurs
\(\sum_k|m_k^{\rm mic}-m_k^{\rm eq}|\) sont
\[
 0.09897380039002812\quad\hbox{et}\quad
 0.9190625715582017,
\]
respectivement. L’énumération parcourt tous les quadruplets d’occupation, calcule leurs multiplicités entières et résout les deux équations de contrainte. Les deux lois se réalisent sur un même écran. Le choix physique entre elles exige une graduation et un Hamiltonien compatibles avec la représentation de la trajectoire. Le paramètre thermique sera typé ultérieurement comme \(\beta=(k_B\Theta)^{-1}\).

\section{Convergence quantifiée vers la distribution de Maxwell--Boltzmann}

La relation entre les deux statistiques quantiques et le régime classique peut
être formulée sans recourir à une approximation informelle. Pour chaque mode, soit
\[
 u_k:=e^{-\beta(E_k^{\rm occ}-\mu)}.
\]
Les occupations par état s’écrivent
\[
 \overline n_k^{\rm FD}=\frac{u_k}{1+u_k},
 \qquad
 \overline n_k^{\rm BE}=\frac{u_k}{1-u_k};
\]
les multiplicités \(g_k\) sont incorporées ensuite au moyen de
\(\langle n_k\rangle=g_k\overline n_k\).

\begin{theorem}[Limite classique commune des statistiques quantiques]
\label{thm:limite-clasico-fd-be-rev2}
Si \(0\leq u_k\leq\delta<1\), alors
\[
 \left|\overline n_k^{\rm FD}-u_k\right|\leq u_k^2,
 \qquad
 \left|\overline n_k^{\rm BE}-u_k\right|
 \leq\frac{u_k^2}{1-\delta}.
\]
Par conséquent, dans le régime dilué \(u_k\to0\),
\[
 \boxed{
 \overline n_k^{\rm FD}\sim
 \overline n_k^{\rm BE}\sim
 e^{-\beta(E_k^{\rm occ}-\mu)}} ,
\]
qui est la distribution de Maxwell--Boltzmann par mode.
\end{theorem}

\begin{proof}
Les identités exactes
\[
 u_k-\frac{u_k}{1+u_k}=\frac{u_k^2}{1+u_k},
 \qquad
 \frac{u_k}{1-u_k}-u_k=\frac{u_k^2}{1-u_k}
\]
fournissent les bornes. Pour \(0<u_k\leq\delta\), en divisant par \(u_k\) et
en prenant \(u_k\downarrow0\), on obtient
l’équivalence asymptotique. Le premier terme d’ordre \(u_k^2\) conserve la
différence physique : l’exclusion diminue l’occupation fermionique, tandis que
l’accumulation augmente l’occupation bosonique.
\end{proof}
